\documentclass{article}
\usepackage{graphicx}
\usepackage{hyperref}

\title{Methodologies in the application of TDA-Mapper}
\author{Clare Kennedy, Darcy, Miller}
\date{April 2026}

\begin{document}

\maketitle

\section{Title: Applications of Topology}

\section{Name, Affiliation, Email}

\begin{itemize}
    \item Lead: George Clare Kennedy, \url{mailto:george-clarekennedy@uiowa.edu}
    \item Co: Jacob Miller, Isabel Darcy
    \item Affiliation: University of Iowa
\end{itemize}

\section{Public Description: 500 Characters}

Topology is easy to think of as a discipline of math with little relevance to the ``real'' world; 
however, the broad idea of finding ``shape'' in data has led to novel insights in various fields, including neuroscience and cancer research. 
In this session, speakers will discuss their use of topological ideas and framekworks in their work, 
as well as dicuss mathematical tools and theory used in topological data analysis.

\section{Sample Speaker List}

\begin{itemize}
    \item Brad Theilman, Sandia National Laboratories
    \item Enrique Alvarado, Iowa State University
\end{itemize}

\section{Description: 1500 Characters}

This session will bring together speakers from many disciplines, under the umbrella of topological data analysis (TDA) and visualization. 
This will include theoretical mathematicians who study algorithms such as Mapper and concepts like persistent homology and dimensionalization reduction, 
as well as researchers who are not mathematicians who have used these ideas successfully in their work. 
There will be ample discussion of software and opportunities to collaborate on larger TDA and visualization libraries, 
an area currently lacking in availability, user friendliness, and standardization. 

\section{MSC}

One of:

\begin{itemize}
    \item 101: Teaching and Learning
    \item 102: Recreational Mathematics
    \item 103: Professional Development and Professional Concerns
    \item 104: Wider Issues
\end{itemize}

\end{document}
